\documentclass[10pt,a4paper]{amsart}
\usepackage{amsmath,amssymb,amsthm,graphicx}
% For commutative diagrams we use Xy-pic
\usepackage[all]{xy}

% We use 2cell for 2-commutative diagrams.
\xyoption{2cell}
% Theorem environments.
%
\theoremstyle{plain}
\newtheorem{theorem}[subsection]{Theorem}
\newtheorem{proposition}[subsection]{Proposition}
\newtheorem{lemma}[subsection]{Lemma}

\theoremstyle{definition}
\newtheorem{definition}[subsection]{Definition}
\newtheorem{example}[subsection]{Example}
\newtheorem{exercise}[subsection]{Exercise}
\newtheorem{situation}[subsection]{Situation}

\theoremstyle{remark}
\newtheorem{remark}[subsection]{Remark}
\newtheorem{remarks}[subsection]{Remarks}

\numberwithin{equation}{subsection}

% Macros
%
\def\lim{\mathop{\mathrm{lim}}\nolimits}
\def\colim{\mathop{\mathrm{colim}}\nolimits}
\def\Spec{\mathop{\mathrm{Spec}}}
\def\Hom{\mathop{\mathrm{Hom}}\nolimits}
\def\Ext{\mathop{\mathrm{Ext}}\nolimits}
\def\SheafHom{\mathop{\mathcal{H}\!\mathit{om}}\nolimits}
\def\SheafExt{\mathop{\mathcal{E}\!\mathit{xt}}\nolimits}
\def\Sch{\mathit{Sch}}
\def\Mor{\mathop{\mathrm{Mor}}\nolimits}
\def\Ob{\mathop{\mathrm{Ob}}\nolimits}
\def\Sh{\mathop{\mathit{Sh}}\nolimits}
\def\NL{\mathop{N\!L}\nolimits}
\def\CH{\mathop{\mathrm{CH}}\nolimits}
\def\proetale{{pro\text{-}\acute{e}tale}}
\def\etale{{\acute{e}tale}}
\def\QCoh{\mathit{QCoh}}
\def\Ker{\mathop{\mathrm{Ker}}}
\def\Im{\mathop{\mathrm{Im}}}
\def\Coker{\mathop{\mathrm{Coker}}}
\def\Coim{\mathop{\mathrm{Coim}}}

% Boxtimes
%
\DeclareMathSymbol{\boxtimes}{\mathbin}{AMSa}{"02}

%
% Macros for moduli stacks/spaces
%
\def\QCohstack{\mathcal{QC}\!\mathit{oh}}
\def\Cohstack{\mathcal{C}\!\mathit{oh}}
\def\Spacesstack{\mathcal{S}\!\mathit{paces}}
\def\Quotfunctor{\mathrm{Quot}}
\def\Hilbfunctor{\mathrm{Hilb}}
\def\Curvesstack{\mathcal{C}\!\mathit{urves}}
\def\Polarizedstack{\mathcal{P}\!\mathit{olarized}}
\def\Complexesstack{\mathcal{C}\!\mathit{omplexes}}
% \Pic is the operator that assigns to X its picard group, usage \Pic(X)
% \Picardstack_{X/B} denotes the Picard stack of X over B
% \Picardfunctor_{X/B} denotes the Picard functor of X over B
\def\Pic{\mathop{\mathrm{Pic}}\nolimits}
\def\Picardstack{\mathcal{P}\!\mathit{ic}}
\def\Picardfunctor{\mathrm{Pic}}
\def\Deformationcategory{\mathcal{D}\!\mathit{ef}}
\usepackage{cleveref}
\crefname{originalreference}{original reference}{original references}
\Crefname{originalreference}{Original reference}{Original references}
\makeatletter
\newlabel{lemma-characterize-finite-locally-constant}{{03RV}{source}{}{sourcepointer-1}{}}
\newlabel{lemma-characterize-finite-locally-constant@cref}{{[originalreference][0][]03RV}{[1][0][]source}{}{}{}}
\makeatother
\crefname{originallemma}{original lemma}{original lemmas}
\Crefname{originallemma}{Original lemma}{Original lemmas}
\makeatletter
\newlabel{pione-lemma-finite-etale-on-proper-over-henselian}{{0A48}{source}{}{sourcepointer-2}{}}
\newlabel{pione-lemma-finite-etale-on-proper-over-henselian@cref}{{[originallemma][0][]0A48}{[1][0][]source}{}{}{}}
\makeatother
\crefname{originallemma}{original lemma}{original lemmas}
\Crefname{originallemma}{Original lemma}{Original lemmas}
\makeatletter
\newlabel{sites-cohomology-lemma-torsors-h1}{{03AJ}{source}{}{sourcepointer-3}{}}
\newlabel{sites-cohomology-lemma-torsors-h1@cref}{{[originallemma][0][]03AJ}{[1][0][]source}{}{}{}}
\makeatother
\title{The first etale cohomology with finite constant coefficients agrees between a proper scheme over a henselian local ring and its closed fibre}
\author{The Stacks Project Authors}
\date{stacks-project:etale-cohomology.tex:local-snapshot-20261011; GNU FDL 1.2 or later}
\begin{document}
\maketitle
\noindent\textit{Editorial note.} All delivered mathematical text is the author's own, copied from the pinned Stacks Project chapter file etale-cohomology.tex; no statement or proof step was rewritten, reordered or removed.\par\medskip
\noindent\textit{Editorial note.} The counted proof is the author's complete argument: the translation of cohomology classes over X into M-torsors represented by schemes finite etale over X, the same translation over the closed fibre, and the identification of the two through the equivalence of categories of finite etale covers.\par\medskip
\noindent\textit{Editorial note.} The label lemma-proper-over-henselian-and-h1 names the counted statement, and the proof environment immediately following it in the source is delivered verbatim as its proof.\par\medskip
\noindent\textit{Editorial note.} The counted statement is carried by the author's own lemma environment in the chapter source.\par\medskip
\noindent\textit{Editorial note.} Pre-open gate: python operations/agent-authoring/tools/precard\_gate.py --source etale-cohomology.tex --label lemma-proper-over-henselian-and-h1 --start 16133 --end 16154 --local-id 3573 returned RESULT=OK (L3573 free, occupancy FREE, public interval clear). Candidate taken from operations/pool-cards/stacks-etale-cohomology-halfB.json.\par\medskip
\noindent\textit{Editorial note.} The statement and proof are copied from the pinned chapter file; the cited inputs are kept as external pointers with their canonical Stacks tags.\par\medskip
\noindent\textit{Allowed standard prerequisites (editorial).}\begin{itemize}
\item The first cohomology with constant finite coefficients classifies torsors. Cohomology on Sites, the lemma identifying the first cohomology group with coefficients in a constant finite group with the set of torsors. Supplies the translation of cohomology classes over X and over X\_0 into torsors, which are finite locally constant sheaves.
\item Representability of a finite locally constant sheaf by a finite etale scheme. Same chapter, the lemma representing a finite locally constant sheaf by a finite etale morphism of schemes. Supplies the translation of the torsors into schemes finite etale over X, respectively over X\_0, with an action of M.
\item Equivalence of categories of finite etale covers of a proper scheme over a henselian local ring and of its closed fibre. Fundamental Groups, the lemma giving an equivalence between finite etale covers of a proper scheme over a henselian local ring and of its closed fibre, compatible with the M-actions. Supplies the identification of the torsor data over X with the torsor data over X\_0, which yields the equality of the cohomology groups.
\end{itemize}\medskip
\ifdefined\hypertarget\hypertarget{sourcepointer-1}{}\fi
\noindent\textit{Source reference.} Original reference 03RV: sources/stacks-etale-cohomology/pinned-original-source.tex, source lines 9994--9994. Prerequisite pointer invoked at proof line(s) 16146; it names a standard fact of the same chapter that the counted argument uses, and every step of the counted proof built on it is retained verbatim in the delivered body.\par\medskip
\ifdefined\hypertarget\hypertarget{sourcepointer-2}{}\fi
\noindent\textit{Source reference.} Original lemma 0A48: sources/stacks-chow/tags.txt, source lines 12766--12766. Prerequisite pointer invoked at proof line(s) 16153; it names a standard fact of a sibling chapter that the counted argument uses, and every step of the counted proof built on it is retained verbatim in the delivered body.\par\medskip
\ifdefined\hypertarget\hypertarget{sourcepointer-3}{}\fi
\noindent\textit{Source reference.} Original lemma 03AJ: sources/stacks-chow/tags.txt, source lines 4312--4312. Prerequisite pointer invoked at proof line(s) 16141; it names a standard fact of a sibling chapter that the counted argument uses, and every step of the counted proof built on it is retained verbatim in the delivered body.\par\medskip
\begin{lemma}
\label{lemma-proper-over-henselian-and-h1}
Let $A$ be a henselian local ring. Let $X$ be a proper scheme over $A$
with closed fibre $X_0$. Let $M$ be a finite abelian group.
Then $H^1_\etale(X, \underline{M}) = H^1_\etale(X_0, \underline{M})$.
\end{lemma}

\begin{proof}
By Cohomology on Sites, Lemma \ref{sites-cohomology-lemma-torsors-h1}
an element of $H^1_\etale(X, \underline{M})$ corresponds to a
$\underline{M}$-torsor $\mathcal{F}$ on $X_\etale$.
Such a torsor is clearly a finite locally constant sheaf.
Hence $\mathcal{F}$ is representable by a scheme $V$ finite
\'etale over $X$, Lemma \ref{lemma-characterize-finite-locally-constant}.
Conversely, a scheme $V$ finite \'etale over $X$ with an $M$-action
which turns it into an $M$-torsor over $X$ gives rise to a cohomology
class. The same translation between cohomology classes over $X_0$ and
torsors finite \'etale over $X_0$ holds. Thus the lemma
is a consequence of the equivalence of categories of
Fundamental Groups, Lemma
\ref{pione-lemma-finite-etale-on-proper-over-henselian}.
\end{proof}

\end{document}
